\import{order/order-type.tex}
\import{set/equinumerosity.tex}

\section{Hartogs's lemma and initial ordinals}

We first record the ordinary transport of a relation along a bijection.  The
carrier arguments keep the set-coded relation explicitly bounded, including
away from its intended domain.

\begin{definition}\label{transported_relation}
    $\transportedRelation{f}{A}{R}{B}
        =\{w\in B\times B\mid
            \exists x,y\in A.\ w=(f(x),f(y))\land x\mathrel{R}y\}$.
\end{definition}

\begin{proposition}\label{transported_relation_binary_on}
    $\transportedRelation{f}{A}{R}{B}$ is a binary relation on $B$.
\end{proposition}
\begin{proof}
    Follows by \obv{transported_relation,subseteq}.
\end{proof}

\begin{proposition}\label{transported_relation_iff}
    Suppose $f\in\bijections{A}{B}$.
    Let $x,y\in A$.
    Then $x\mathrel{R}y$ iff
        $f(x)\mathrel{\transportedRelation{f}{A}{R}{B}}f(y)$.
\end{proposition}
\begin{proof}
    $f\in\funs{A}{B}$ and $f$ is a function and $\dom{f}=A$
        by \obv{bijections_to_funs,bijection_is_function,bijections_dom}.
    $f$ is injective by \obv{bijections}.
    $f(x),f(y)\in B$ by \obv{funs_type_apply}.
    Show that if $x\mathrel{R}y$, then
        $f(x)\mathrel{\transportedRelation{f}{A}{R}{B}}f(y)$.
    \begin{subproof}
        Assume $x\mathrel{R}y$.
        Follows by
            \obv{transported_relation,times_tuple_intro,pair_eq_iff}.
    \end{subproof}
    Show that if
        $f(x)\mathrel{\transportedRelation{f}{A}{R}{B}}f(y)$,
        then $x\mathrel{R}y$.
    \begin{subproof}
        Assume
            $f(x)\mathrel{\transportedRelation{f}{A}{R}{B}}f(y)$.
        Take $u,v\in A$ such that
            $(f(x),f(y))=(f(u),f(v))$ and $u\mathrel{R}v$
            by \obv{transported_relation}.
        $f(x)=f(u)$ and $f(y)=f(v)$ by \obv{pair_eq_iff}.
        Hence $x=u$ and $y=v$ by \auto{injective_function}.
    \end{subproof}
    Follows by \auto.
\end{proof}

\begin{theorem}\label{transported_relation_isomorphism}
    Suppose $f\in\bijections{A}{B}$.
    Suppose $R$ is a binary relation on $A$.
    Then
        $\orderIsomorphism{f}{A}{R}{B}
            {\transportedRelation{f}{A}{R}{B}}$.
\end{theorem}
\begin{proof}
    $\transportedRelation{f}{A}{R}{B}$ is a binary relation on $B$
        by \obv{transported_relation_binary_on}.
    For all $x,y\in A$ we have $x\mathrel{R}y$ iff
        $f(x)\mathrel{\transportedRelation{f}{A}{R}{B}}f(y)$
        by \obv{transported_relation_iff}.
    Follows by \obv{order_isomorphism}.
\end{proof}

\begin{proposition}\label{unions_of_family_of_ordinals_is_ordinal}
    Suppose every element of $F$ is an ordinal.
    Then $\unions{F}$ is an ordinal.
\end{proposition}
\begin{proof}
    Show that $\unions{F}$ is \in-transitive.
    \begin{subproof}
        Show that for all $x,y$ such that $x\in y\in\unions{F}$
            we have $x\in\unions{F}$.
        \begin{subproof}
            Fix $x,y$.
            Assume $x\in y\in\unions{F}$.
            Take $\alpha\in F$ such that $y\in\alpha$
                by \obv{unions_iff}.
            Hence $x\in\alpha$ by \obv{ordinal_transitivity}.
            Follows by \auto{unions_iff}.
        \end{subproof}
        Follows by \obv{transitiveset}.
    \end{subproof}
    Show that every element of $\unions{F}$ is \in-transitive.
    \begin{subproof}
        Fix $x\in\unions{F}$.
        Take $\alpha\in F$ such that $x\in\alpha$
            by \obv{unions_iff}.
        $x$ is an ordinal by \obv{elem_of_ordinal_is_ordinal}.
        Follows by \obv{ordinal}.
    \end{subproof}
    Follows by \obv{ordinal}.
\end{proof}

The following family is a set: both the carrier and its relation are selected
from powersets determined by $X$.  Thus its order types can be collected by
ordinary replacement without quantifying over all well-orders as a class.

\begin{definition}\label{well_order_codes}
    $\wellOrderCodes{X}
        =\{p\in\pow{X}\times\pow{X\times X}\mid
            \text{$\snd{p}$ is a strict well order on $\fst{p}$}\}$.
\end{definition}

\begin{definition}\label{well_order_types}
    $\wellOrderTypes{X}
        =\{\orderType{\fst{p}}{\snd{p}}\mid
            p\in\wellOrderCodes{X}\}$.
\end{definition}

\begin{definition}\label{hartogs_successors}
    $\hartogsSuccessors{X}
        =\{\suc{\alpha}\mid\alpha\in\wellOrderTypes{X}\}$.
\end{definition}

\begin{definition}\label{hartogs}
    $\hartogs{X}=\unions{\hartogsSuccessors{X}}$.
\end{definition}

\begin{proposition}\label{well_order_type_family_is_ordinal}
    Every element of $\wellOrderTypes{X}$ is an ordinal.
\end{proposition}
\begin{proof}
    Fix $\alpha\in\wellOrderTypes{X}$.
    Take $p\in\wellOrderCodes{X}$ such that
        $\alpha=\orderType{\fst{p}}{\snd{p}}$
        by \obv{well_order_types}.
    $\snd{p}$ is a strict well order on $\fst{p}$
        by \obv{well_order_codes}.
    Follows by \obv{order_type_is_ordinal}.
\end{proof}

\begin{proposition}\label{hartogs_successor_family_is_ordinal}
    Every element of $\hartogsSuccessors{X}$ is an ordinal.
\end{proposition}
\begin{proof}
    Fix $\beta\in\hartogsSuccessors{X}$.
    Take $\alpha\in\wellOrderTypes{X}$ such that $\beta=\suc{\alpha}$
        by \obv{hartogs_successors}.
    $\alpha$ is an ordinal by \obv{well_order_type_family_is_ordinal}.
    Follows by \obv{suc_ordinal}.
\end{proof}

\begin{theorem}\label{hartogs_is_ordinal}
    $\hartogs{X}$ is an ordinal.
\end{theorem}
\begin{proof}
    Follows by
        \auto{hartogs,unions_of_family_of_ordinals_is_ordinal,hartogs_successor_family_is_ordinal}.
\end{proof}

\begin{proposition}\label{order_type_belongs_to_hartogs}
    Suppose $A\subseteq X$.
    Suppose $R$ is a strict well order on $A$.
    Then $\orderType{A}{R}\in\hartogs{X}$.
\end{proposition}
\begin{proof}
    $A\in\pow{X}$ by \obv{powerset_intro}.
    $R\subseteq A\times A$
        by \auto{strict_well_order,strict_linear_order}.
    $A\times A\subseteq X\times X$
        by \auto{times_subseteq_left,times_subseteq_right,subseteq_transitive}.
    Hence $R\subseteq X\times X$ by \obv{subseteq_transitive}.
    Hence $R\in\pow{X\times X}$ by \obv{powerset_intro}.
    Hence $(A,R)\in\pow{X}\times\pow{X\times X}$
        by \obv{times_tuple_intro}.
    Hence $(A,R)\in\wellOrderCodes{X}$
        by \obv{well_order_codes,fst_eq,snd_eq}.
    Hence $\orderType{A}{R}\in\wellOrderTypes{X}$
        by \obv{well_order_types,fst_eq,snd_eq}.
    Hence $\suc{\orderType{A}{R}}\in\hartogsSuccessors{X}$
        by \obv{hartogs_successors}.
    $\orderType{A}{R}\in\suc{\orderType{A}{R}}$
        by \obv{suc_intro_self}.
    Hence $\orderType{A}{R}\in\unions{\hartogsSuccessors{X}}$
        by \auto{unions_iff}.
    Follows by \obv{hartogs}.
\end{proof}

\begin{theorem}[Hartogs]\label{hartogs_theorem}
    There exists no $f$ such that $f\in\inj{\hartogs{X}}{X}$.
\end{theorem}
\begin{proof}
    Assume there exists $f$ such that $f\in\inj{\hartogs{X}}{X}$.
    Take $f$ such that $f\in\inj{\hartogs{X}}{X}$.
    Let $H=\hartogs{X}$.
    Let $Y=\ran{f}$.
    $f\in\funs{H}{X}$ by \obv{inj}.
    Hence $f$ is a function and $\dom{f}=H$
        by \obv{funs_is_function,funs}.
    $f$ is injective by \auto{inj,injective_function}.
    $Y\subseteq X$ by \auto{funs_ran}.
    $f\in\surj{H}{Y}$ by \auto{fun_to_surj}.
    Hence $f\in\bijections{H}{Y}$ by \obv{bijections}.
    $H$ is an ordinal by \obv{hartogs_is_ordinal}.
    Hence $\memrel{H}$ is a strict well order on $H$
        by \obv{ordinal_membership_is_strict_well_order}.
    Let $S=\transportedRelation{f}{H}{\memrel{H}}{Y}$.
    $\orderIsomorphism{f}{H}{\memrel{H}}{Y}{S}$
        by
        \auto{transported_relation_isomorphism,strict_well_order,strict_linear_order}.
    Hence $S$ is a strict well order on $Y$
        by \obv{order_isomorphism_preserves_strict_well_order}.
    $\orderIsomorphism{\converse{f}}{Y}{S}{H}{\memrel{H}}$
        by \obv{order_isomorphism_inverse}.
    Hence $\orderIsomorphic{Y}{S}{H}{\memrel{H}}$
        by \obv{order_isomorphic}.
    Hence $H=\orderType{Y}{S}$
        by \obv{ordinal_order_type_unique}.
    $\orderType{Y}{S}\in\hartogs{X}$
        by \obv{order_type_belongs_to_hartogs}.
    Hence $H\in H$.
    Contradiction by \obv{in_irrefl}.
\end{proof}

\begin{theorem}\label{below_hartogs_injection}
    Let $\alpha\in\hartogs{X}$.
    Then there exists $f\in\inj{\alpha}{X}$.
\end{theorem}
\begin{proof}
    $\hartogs{X}$ is an ordinal by \obv{hartogs_is_ordinal}.
    Hence $\alpha$ is an ordinal by \obv{elem_of_ordinal_is_ordinal}.
    $\alpha\in\unions{\hartogsSuccessors{X}}$ by \obv{hartogs}.
    Take $\gamma\in\hartogsSuccessors{X}$ such that $\alpha\in\gamma$
        by \obv{unions_iff}.
    Take $\beta\in\wellOrderTypes{X}$ such that $\gamma=\suc{\beta}$
        by \obv{hartogs_successors}.
    $\beta$ is an ordinal by \obv{well_order_type_family_is_ordinal}.
    $\alpha=\beta$ or $\alpha\in\beta$ by \obv{suc_iff}.
    Hence $\alpha\subseteq\beta$
        by \obv{subseteq_refl,ordinal_elem_implies_subseteq}.
    Take $p\in\wellOrderCodes{X}$ such that
        $\beta=\orderType{\fst{p}}{\snd{p}}$
        by \obv{well_order_types}.
    Let $A=\fst{p}$.
    Let $R=\snd{p}$.
    $p\in\pow{X}\times\pow{X\times X}$ and
        $R$ is a strict well order on $A$ by \auto{well_order_codes}.
    Hence $A\in\pow{X}$ by \auto{fst_type}.
    Hence $A\subseteq X$ by \obv{powerset_elim}.
    $\orderIsomorphism{\orderCollapse{A}{R}}
        {A}{R}{\beta}{\memrel{\beta}}$
        by \auto{order_collapse_is_isomorphism}.
    Hence $\orderIsomorphism{\converse{\orderCollapse{A}{R}}}
        {\beta}{\memrel{\beta}}{A}{R}$
        by \obv{order_isomorphism_inverse}.
    Let $g=\converse{\orderCollapse{A}{R}}$.
    $g\in\bijections{\beta}{A}$ by \auto{order_isomorphism}.
    Hence $g\in\funs{\beta}{A}$ and $g$ is injective
        by \obv{bijections,bijections_to_funs}.
    Hence $g\in\funs{\beta}{X}$ by \obv{funs_weaken_codom}.
    Hence $g\in\inj{\beta}{X}$
        by \auto{inj,funs_is_function,funs,injective_function}.
    $\restrl{g}{\alpha}\in\inj{\alpha}{X}$ by \obv{inj_restrl}.
\end{proof}

An initial ordinal is an ordinal which is not equinumerous with any smaller
ordinal.  Cardinals will be introduced as a synonym for these objects in the
next layer.

\begin{definition}\label{initial_ordinal}
    $\kappa$ is an initial ordinal iff $\kappa$ is an ordinal and
        for every $\alpha\in\kappa$ we have
            $\lnot(\alpha\approx\kappa)$.
\end{definition}

\begin{theorem}\label{initial_ordinal_le_equinumerous_ordinal}
    Suppose $\kappa$ is an initial ordinal.
    Let $\alpha$ be an ordinal.
    Suppose $\alpha\approx\kappa$.
    Then $\kappa\subseteq\alpha$.
\end{theorem}
\begin{proof}
    Suppose not.
    Hence $\alpha\subseteq\kappa$
        by \obv{ordinal_subseteq_or_subseteq}.
    Hence $\alpha\in\kappa$ or $\alpha=\kappa$
        by \obv{ordinal_subseteq_implies_elem_or_eq}.
    Hence $\alpha\in\kappa$ by \obv{subseteq_refl}.
    $\lnot(\alpha\approx\kappa)$ by \obv{initial_ordinal}.
    Contradiction.
\end{proof}

\begin{proposition}\label{initial_ordinal_characterization}
    $\kappa$ is an initial ordinal iff $\kappa$ is an ordinal and
        for every ordinal $\alpha$ such that $\alpha\approx\kappa$
            we have $\kappa\subseteq\alpha$.
\end{proposition}
\begin{proof}
    Show that if $\kappa$ is an initial ordinal, then $\kappa$ is an
        ordinal and for every ordinal $\alpha$ such that
        $\alpha\approx\kappa$ we have $\kappa\subseteq\alpha$.
    \begin{subproof}
        Assume $\kappa$ is an initial ordinal.
        For every ordinal $\alpha$ such that $\alpha\approx\kappa$
            we have $\kappa\subseteq\alpha$
            by \obv{initial_ordinal_le_equinumerous_ordinal}.
    \end{subproof}
    Show that if $\kappa$ is an ordinal and for every ordinal $\alpha$
        such that $\alpha\approx\kappa$ we have
        $\kappa\subseteq\alpha$, then $\kappa$ is an initial ordinal.
    \begin{subproof}
        Assume $\kappa$ is an ordinal and for every ordinal $\alpha$
            such that $\alpha\approx\kappa$ we have
            $\kappa\subseteq\alpha$.
        Show that for every $\alpha\in\kappa$ we have
            $\lnot(\alpha\approx\kappa)$.
        \begin{subproof}
            Fix $\alpha\in\kappa$.
            $\alpha$ is an ordinal by \obv{elem_of_ordinal_is_ordinal}.
            Suppose $\alpha\approx\kappa$.
            Hence $\kappa\subseteq\alpha$.
            Hence $\alpha\in\alpha$ by \obv{subseteq}.
            Contradiction by \obv{in_irrefl}.
        \end{subproof}
        Follows by \obv{initial_ordinal}.
    \end{subproof}
    Follows by \auto.
\end{proof}

\begin{proposition}\label{equinumerous_initial_ordinals_are_equal}
    Let $\kappa,\lambda$ be initial ordinals.
    Suppose $\kappa\approx\lambda$.
    Then $\kappa=\lambda$.
\end{proposition}
\begin{proof}
    $\kappa\subseteq\lambda$
        by \obv{initial_ordinal_le_equinumerous_ordinal,equinum_sym}.
    $\lambda\subseteq\kappa$
        by \obv{initial_ordinal_le_equinumerous_ordinal}.
    Follows by \obv{subseteq_antisymmetric}.
\end{proof}
